% Generated draft: requires mathematical and source review.
\documentclass[11pt]{article}
\usepackage[margin=25mm]{geometry}
\usepackage{fontspec}
\setmainfont{Latin Modern Roman}
\usepackage{amsmath,amssymb,mathtools,unicode-math}
\setmathfont{Latin Modern Math}
\usepackage{enumitem,needspace,xurl,hyperref}
\setlength{\parindent}{0pt}
\setlength{\parskip}{5pt}
\setlength{\emergencystretch}{3em}
\newcommand{\N}{\mathbb{N}}
\newcommand{\Z}{\mathbb{Z}}
\newcommand{\Q}{\mathbb{Q}}
\newcommand{\R}{\mathbb{R}}
\newcommand{\C}{\mathbb{C}}
\newcommand{\F}{\mathbb{F}}
\newcommand{\A}{\mathbb{A}}
\newcommand{\PP}{\mathbb{P}}
\newcommand{\E}{\mathbb{E}}
\newcommand{\eps}{\varepsilon}
\newcommand{\dd}{\,\mathrm{d}}
\begin{document}
\section*{UnsolvedMath ID: 20002886}
\textbf{Source problem code:} catalogue entry
\subsection{Definitions and mathematical statement}
Interpret $X,Y$ as connected compact Riemann surfaces without boundary. Let $a_1,\ldots,a_s\in X$ be distinct marked points and let $m_1,\ldots,m_s\in\Z_{>0}$. The source asks for a nonconstant holomorphic map $f:X\to Y$ with ramification divisor exactly
\[
\operatorname{Ram}(f)=\sum_{j=1}^s m_j[a_j].
\]
Equivalently, there are centered local source and target coordinates at $a_j$ and $f(a_j)$ in which $f$ is $\zeta\mapsto\zeta^{m_j+1}$, and $f$ has no critical points outside the marked set. The target coordinates may depend on $j$; no common target value is prescribed. This repairs the extracted superscript in ``$z^{m_j+1}$''.

For degree $d=\deg f\geq1$, and genera $g_X,g_Y$, the supplied necessary conditions are
\[
\sum_{j=1}^s m_j=2g_X-2-d(2g_Y-2),\qquad m_j+1\leq d,
\]
and, for every $y\in Y$,
\[
\sum_{\{j:f(a_j)=y\}}(m_j+1)\leq d.
\]
The last condition is a fiberwise constraint; the unknown critical values are part of the map to be found. These numerical conditions are not asserted sufficient for arbitrary fixed complex structures and marked points.

The aid gives a complete restricted classification when $X=Y=\PP^1$ and exactly two critical points $a_1,a_2$ are prescribed. Existence holds if and only if $m_1=m_2$. In that case $d=m_1+1=m_2+1$, the critical values are distinct, and the maps have the form
\[
f=B\circ(\zeta\mapsto\zeta^d)\circ A,
\]
where $A,B$ are M\"obius transformations and $A(a_1)=0$, $A(a_2)=\infty$ (with the order of the two points interchangeable). This classification concerns marked source points but unprescribed target critical values.

The source next proposes an additional hypothesis: a globally smooth $\phi:X\to\R$ with $\Delta\phi\geq0$ everywhere, $\Delta\phi>0$ off the $a_j$, and a local relation between $\Delta\phi$ and $|\zeta|^{2m_j}$ near $a_j$. The comparison symbol in that last relation is missing in the extraction, so equality, comparability, or another condition cannot be chosen silently. Nevertheless the preceding global inequalities alone are inconsistent on a compact boundaryless $X$. For any smooth metric and its Laplace--Beltrami operator $\Delta=\operatorname{div}\operatorname{grad}$,
\[
\int_X\Delta\phi\,\dd V=0,
\qquad \Delta\phi\geq0\ \Longrightarrow\ \Delta\phi\equiv0,
\]
contradicting strict positivity on the complement of the finite marked set. Thus the proposed extra-potential hypothesis is vacuous in the literal compact setting, irrespective of the lost local comparison symbol. Replacing compactness, allowing singular potentials, or adding boundary data would produce different problems and is not assumed.
\subsection{Short English statement}
The first question prescribes every ramification point and its order for a holomorphic map between fixed compact Riemann surfaces. The supplied genus, local-degree, and fiber constraints are necessary, and the two-critical-point sphere case reduces exactly to power maps with equal multiplicities. The proposed smooth subharmonic potential cannot exist on a compact surface without boundary, so that additional hypothesis cannot supply a nonvacuous existence criterion as printed.
\subsection{Sources}
\begin{itemize}
\item Local registry: permanent catalogue entry, exact source code catalogue entry; \texttt{aim-several-complex-variables-notes.json}, \texttt{notes[43]}.
\item AIM workshop list \emph{Complexity of mappings in CR geometry}, Problem 4, ``Mappings between Riemann surfaces'', proposed by Dror Varolin, \url{https://aimath.org/WWN/complexitycrmap/complexitycrmap.pdf}. The original remarks and literature field are empty; no full authoritative publication reference or proposal year is supplied.
\item Verification field reports corrupted superscripts and a lost comparison symbol, but ends before giving its promised recovered formulation.
\item Supplied research attempt 1, registry updated 14 August 2026: partial result, low novelty confidence, machine-generated and not peer reviewed. It supplies the Riemann--Hurwitz and fiber constraints, two-point spherical classification, and compact-Laplacian obstruction without a full external citation.
\end{itemize}
\end{document}
